% GENERATED FILE. Edit canonical lessons, then run npm run generate:books.
% canonical-source-hash: fnv1a64:3c60ada9a99502b2
% canonical-lessons: SA-C220-tula, SA-C220-manadandah, SA-C220-dandah, SA-C220-yastih, SA-C220-dhanuh, SA-R220-first-pass-words-from-chapters-212-215, SA-R220-second-pass-words-from-chapters-216-220

\chapter{Balance, Measuring rod, Stick, Staff, Archer's bow}
\label{ch:sa-balance-measuring-rod-stick-staff-archer-s-bow}
\hlchaptermodality{220}

\begin{chapteropening}
\textbf{By the end of this chapter:} \emph{I can say a balance, a measuring rod, a stick, a staff and an archer's bow.}

Complete the last lesson of chapter 220: I can say a balance, a measuring rod, a stick, a staff and an archer's bow.
\end{chapteropening}

\section[\texorpdfstring{\sk{तुला} (tulā)}{तुला (tulā)}]{\sk{तुला} (tulā) — a balance, scales}

\label{lesson:SA-C220-tula}



\begin{quote}
Before the new one: say the Sanskrit for an apple, then the Sanskrit for flattened rice.
\end{quote}

\subsection*{You'll want to know: \sk{तुला}}
\textbf{\sk{तुला}} — \emph{tulā} — \textquotedblleft{}a balance, scales\textquotedblright{}.

\begin{grammarlens}[title={What you've built}]
One more word for this chapter.
\end{grammarlens}

\subsection*{Your turn}
\begin{itemize}
\raggedright
  \item \emph{Say it:} \emph{tulā}
  \item \emph{Say it:} \emph{tulā}, once more
  \item \emph{Say it:} say \emph{tulā}
  \item \emph{Recall:} say the Sanskrit for an apple, then the Sanskrit for flattened rice, then say \emph{tulā} again
\end{itemize}

\begin{tcolorbox}[breakable,colback=teal!4,colframe=teal!35!black,title={Before you move on}]
What does \sk{तुला} mean? (\textquotedblleft{}A balance, scales\textquotedblright{}.) Say it once more.
\end{tcolorbox}
\section[\texorpdfstring{\sk{मानदण्डः} (mānadaṇḍaḥ)}{मानदण्डः (mānadaṇḍaḥ)}]{\sk{मानदण्डः} (mānadaṇḍaḥ) — a measuring rod}

\label{lesson:SA-C220-manadandah}



\begin{quote}
Before the new one: say the Sanskrit for flattened rice, then the Sanskrit for a balance.
\end{quote}

\subsection*{You'll want to know: \sk{मानदण्डः}}
\textbf{\sk{मानदण्डः}} — \emph{mānadaṇḍaḥ} — \textquotedblleft{}a measuring rod\textquotedblright{}.

\begin{grammarlens}[title={What you've built}]
One more word for this chapter.
\end{grammarlens}

\subsection*{Your turn}
\begin{itemize}
\raggedright
  \item \emph{Say it:} \emph{mānadaṇḍaḥ}
  \item \emph{Say it:} \emph{mānadaṇḍaḥ}, once more
  \item \emph{Say it:} say \emph{mānadaṇḍaḥ}
  \item \emph{Recall:} say the Sanskrit for flattened rice, then the Sanskrit for a balance, then say \emph{mānadaṇḍaḥ} again
\end{itemize}

\begin{tcolorbox}[breakable,colback=teal!4,colframe=teal!35!black,title={Before you move on}]
What does \sk{मानदण्डः} mean? (\textquotedblleft{}A measuring rod\textquotedblright{}.) Say it once more.
\end{tcolorbox}
\section[\texorpdfstring{\sk{दण्डः} (daṇḍaḥ)}{दण्डः (daṇḍaḥ)}]{\sk{दण्डः} (daṇḍaḥ) — a stick}

\label{lesson:SA-C220-dandah}



\begin{quote}
Before the new one: say the Sanskrit for a balance, then the Sanskrit for a measuring rod.
\end{quote}

\subsection*{You'll want to know: \sk{दण्डः}}
\textbf{\sk{दण्डः}} — \emph{daṇḍaḥ} — \textquotedblleft{}a stick\textquotedblright{}.

\begin{grammarlens}[title={What you've built}]
One more word for this chapter.
\end{grammarlens}

\subsection*{Your turn}
\begin{itemize}
\raggedright
  \item \emph{Say it:} \emph{daṇḍaḥ}
  \item \emph{Say it:} \emph{daṇḍaḥ}, once more
  \item \emph{Say it:} say \emph{daṇḍaḥ}
  \item \emph{Recall:} say the Sanskrit for a balance, then the Sanskrit for a measuring rod, then say \emph{daṇḍaḥ} again
\end{itemize}

\begin{tcolorbox}[breakable,colback=teal!4,colframe=teal!35!black,title={Before you move on}]
What does \sk{दण्डः} mean? (\textquotedblleft{}A stick\textquotedblright{}.) Say it once more.
\end{tcolorbox}
\section[\texorpdfstring{\sk{यष्टिः} (yaṣṭiḥ)}{यष्टिः (yaṣṭiḥ)}]{\sk{यष्टिः} (yaṣṭiḥ) — a staff}

\label{lesson:SA-C220-yastih}



\begin{quote}
Before the new one: say the Sanskrit for a measuring rod, then the Sanskrit for a stick.
\end{quote}

\subsection*{You'll want to know: \sk{यष्टिः}}
\textbf{\sk{यष्टिः}} — \emph{yaṣṭiḥ} — \textquotedblleft{}a staff\textquotedblright{}.

\begin{grammarlens}[title={What you've built}]
One more word for this chapter.
\end{grammarlens}

\subsection*{Your turn}
\begin{itemize}
\raggedright
  \item \emph{Say it:} \emph{yaṣṭiḥ}
  \item \emph{Say it:} \emph{yaṣṭiḥ}, once more
  \item \emph{Say it:} say \emph{yaṣṭiḥ}
  \item \emph{Recall:} say the Sanskrit for a measuring rod, then the Sanskrit for a stick, then say \emph{yaṣṭiḥ} again
\end{itemize}

\begin{tcolorbox}[breakable,colback=teal!4,colframe=teal!35!black,title={Before you move on}]
What does \sk{यष्टिः} mean? (\textquotedblleft{}A staff\textquotedblright{}.) Say it once more.
\end{tcolorbox}
\section[\texorpdfstring{\sk{धनुः} (dhanuḥ)}{धनुः (dhanuḥ)}]{\sk{धनुः} (dhanuḥ) — an archer's bow}

\label{lesson:SA-C220-dhanuh}



\begin{quote}
Before the new one: say the Sanskrit for a stick, then the Sanskrit for a staff.
\end{quote}

\subsection*{You'll want to know: \sk{धनुः}}
\textbf{\sk{धनुः}} — \emph{dhanuḥ} — \textquotedblleft{}an archer's bow\textquotedblright{}.

\begin{grammarlens}[title={What you've built}]
One more word for this chapter.
\end{grammarlens}

\subsection*{Your turn}
\begin{itemize}
\raggedright
  \item \emph{Say it:} \emph{dhanuḥ}
  \item \emph{Say it:} \emph{dhanuḥ}, once more
  \item \emph{Say it:} say \emph{dhanuḥ}
  \item \emph{Recall:} say the Sanskrit for a stick, then the Sanskrit for a staff, then say \emph{dhanuḥ} again
\end{itemize}

\begin{tcolorbox}[breakable,colback=teal!4,colframe=teal!35!black,title={Before you move on}]
What does \sk{धनुः} mean? (\textquotedblleft{}An archer's bow\textquotedblright{}.) Say it once more.
\end{tcolorbox}
\section[(dialogue)]{Words from chapters 212-215 — a first pass}

\label{lesson:SA-R220-first-pass-words-from-chapters-212-215}



\begin{quote}
Say the words for a staff and an archer's bow.
\end{quote}

\subsection*{Your turn}
Say these aloud:

\begin{itemize}
\raggedright
  \item intelligence, the mind, the self, a dream, satisfaction
  \item at once, instantly, formerly, afterwards, as long as
  \item exactly, indeed, like, without, towards
  \item moving, slow, sharp, blunt, valuable
\end{itemize}

\begin{tcolorbox}[breakable,colback=teal!4,colframe=teal!35!black,title={Before you move on}]
Intelligence? (\textbf{\sk{बुद्धिः}}, \emph{buddhiḥ}.) The mind? (\textbf{\sk{मनः}}, \emph{manaḥ}.) The self? (\textbf{\sk{आत्मा}}, \emph{ātmā}.) A dream? (\textbf{\sk{स्वप्नः}}, \emph{svapnaḥ}.) Satisfaction? (\textbf{\sk{तृप्तिः}}, \emph{tṛptiḥ}.) At once? (\textbf{\sk{सद्यः}}, \emph{sadyaḥ}.) Instantly? (\textbf{\sk{तत्क्षणम्}}, \emph{tatkṣaṇam}.) Formerly? (\textbf{\sk{प्राक्}}, \emph{prāk}.) Afterwards? (\textbf{\sk{अनन्तरम्}}, \emph{anantaram}.) As long as? (\textbf{\sk{यावत्}}, \emph{yāvat}.) Exactly? (\textbf{\sk{एव}}, \emph{eva}.) Indeed? (\textbf{\sk{खलु}}, \emph{khalu}.) Like? (\textbf{\sk{इव}}, \emph{iva}.) Without? (\textbf{\sk{विना}}, \emph{vinā}.) Towards? (\textbf{\sk{प्रति}}, \emph{prati}.) Moving? (\textbf{\sk{चलः}}, \emph{calaḥ}.) Slow? (\textbf{\sk{मन्दः}}, \emph{mandaḥ}.) Sharp? (\textbf{\sk{तीक्ष्णः}}, \emph{tīkṣṇaḥ}.) Blunt? (\textbf{\sk{कुण्ठः}}, \emph{kuṇṭhaḥ}.) Valuable? (\textbf{\sk{मूल्यवान्}}, \emph{mūlyavān}.)
\end{tcolorbox}
\section[(dialogue)]{Words from chapters 216-220 — a second pass}

\label{lesson:SA-R220-second-pass-words-from-chapters-216-220}



\begin{quote}
Say the words for able and expert.
\end{quote}

\subsection*{Your turn}
Say these aloud:

\begin{itemize}
\raggedright
  \item able, expert, ignorant, wise, intelligent
  \item red, reddish, tawny, equal, uneven
  \item a broom, an axe, a spade, a plough, a sickle
  \item a bael fruit, an orange, a wood apple, an apple, flattened rice
  \item a balance, a measuring rod, a stick, a staff, an archer's bow
\end{itemize}

\begin{tcolorbox}[breakable,colback=teal!4,colframe=teal!35!black,title={Before you move on}]
Able? (\textbf{\sk{कुशलः}}, \emph{kuśalaḥ}.) Expert? (\textbf{\sk{निपुणः}}, \emph{nipuṇaḥ}.) Ignorant? (\textbf{\sk{अज्ञः}}, \emph{ajñaḥ}.) Wise? (\textbf{\sk{विज्ञः}}, \emph{vijñaḥ}.) Intelligent? (\textbf{\sk{बुद्धिमान्}}, \emph{buddhimān}.) Red? (\textbf{\sk{रक्तः}}, \emph{raktaḥ}.) Reddish? (\textbf{\sk{अरुणः}}, \emph{aruṇaḥ}.) Tawny? (\textbf{\sk{कपिलः}}, \emph{kapilaḥ}.) Equal? (\textbf{\sk{समः}}, \emph{samaḥ}.) Uneven? (\textbf{\sk{विषमः}}, \emph{viṣamaḥ}.) A broom? (\textbf{\sk{सम्मार्जनी}}, \emph{sammārjanī}.) An axe? (\textbf{\sk{कुठारः}}, \emph{kuṭhāraḥ}.) A spade? (\textbf{\sk{खनित्रम्}}, \emph{khanitram}.) A plough? (\textbf{\sk{हलम्}}, \emph{halam}.) A sickle? (\textbf{\sk{दात्रम्}}, \emph{dātram}.) A bael fruit? (\textbf{\sk{बिल्वम्}}, \emph{bilvam}.) An orange? (\textbf{\sk{नारंगम्}}, \emph{nāraṅgam}.) A wood apple? (\textbf{\sk{कपित्थः}}, \emph{kapitthaḥ}.) An apple? (\textbf{\sk{सेवफलम्}}, \emph{sevaphalam}.) Flattened rice? (\textbf{\sk{चिपिटकः}}, \emph{cipiṭakaḥ}.) A balance? (\textbf{\sk{तुला}}, \emph{tulā}.) A measuring rod? (\textbf{\sk{मानदण्डः}}, \emph{mānadaṇḍaḥ}.) A stick? (\textbf{\sk{दण्डः}}, \emph{daṇḍaḥ}.) A staff? (\textbf{\sk{यष्टिः}}, \emph{yaṣṭiḥ}.) An archer's bow? (\textbf{\sk{धनुः}}, \emph{dhanuḥ}.)
\end{tcolorbox}
